\documentclass[11pt,a4paper]{article}
\usepackage[T1]{fontenc}
\usepackage{lmodern,microtype}
\usepackage[a4paper,margin=25mm]{geometry}
\usepackage{amsmath,amssymb,booktabs,array,tabularx,longtable}
\usepackage[dvipsnames]{xcolor}
\usepackage[colorlinks=true,linkcolor=MidnightBlue,urlcolor=MidnightBlue]{hyperref}
\usepackage{fancyhdr}
\pagestyle{fancy}
\fancyhf{}
\fancyhead[L]{\small Rank-five fusion categories}
\fancyhead[R]{\small Detailed revision letter}
\fancyfoot[C]{\thepage}
\setlength{\headheight}{14pt}
\setlength{\emergencystretch}{2em}
\setlength{\parskip}{3pt}
\newcommand{\FPdim}{\operatorname{FPdim}}
\newcommand{\Tr}{\operatorname{Tr}}
\hypersetup{pdftitle={Rank-five classification: second detailed audit and revision letter},pdfauthor={}}
\title{Rank-five classification\\[4pt]\large Second detailed audit and revision letter}
\author{}
\date{24 November 2026}
\begin{document}
\maketitle

\section{Outcome and exact scope}

This letter records the new review of the supplied
\path{Rank5_General_Sharp.zip}, whose manuscript has 69 pages.
It concerns changes relative to that archive, including changes to its
verification programs. It supersedes the earlier review as an account of
the \emph{current} audit; the supplied \path{Review_Report.tex} and
\path{Review_Report.pdf} are preserved as historical documents.

The review found two statements needing a substantive qualification,
one missing conductor hypothesis, an insufficiently justified numerical
upper endpoint, and several reproducibility and exposition defects.
The revised manuscript repairs them. The main conclusions are retained:
\[
\begin{gathered}
\text{every complex rank-five fusion category has a pseudo-unitary Galois
conjugate}\\[-2pt]
\text{and admits a spherical structure},\\
\mu\le4,\qquad \FPdim(\mathcal C)\le48+24\sqrt3,\\
\text{exactly seventeen based rings, each with a unitary realization}.
\end{gathered}
\]
No counterexample to these conclusions or unrepaired gap in the reviewed
dependency chain was identified. This is a mathematical and computational
audit, not a proof-assistant certification or a guarantee that no error
can remain.

The review read all seventeen sections, examined the ordinary/signed
induction and Galois descent arguments, the invertible-action and Morita
arguments, the character-order and twist-support lemmas, and the
field-by-field arithmetic reductions. It also read the eighteen structural
programs, the three census parametrizations, the generators, exclusion
criteria, witness format, model registry and documentation. The original
checksum manifest and the complete original verification both passed.
The revised verification passed again after the changes below.

Change identifiers M01--M18 and C01--C04 provide a complete ledger of
manuscript and program edits. Section~\ref{sec:package} records the
remaining documentation and delivery changes. The accompanying
\path{revision.patch} gives the exact textual differences in existing
source files.

\section{Statements corrected}

\subsection{Pivotality and the existence of spherical structures: M01, M03}

\textbf{Location: Proposition 3.2 and the end of its proof.}
The old statement, ``every pivotal rank-five category is spherical'',
was ambiguous about whether the chosen pivotal structure was meant.
The proof established existence of a spherical structure in the graded
case. It did not show that every pivotal structure on a graded category
is spherical.

That stronger interpretation is false. On $\mathrm{Vec}_{C_5}$, twisting
the canonical pivotal structure by a character taking a generator to a
primitive fifth root $\zeta_5$ gives left and right simple dimensions
$\zeta_5$ and $\zeta_5^{-1}$, in one order. They differ.

The revised statement separates the valid assertions:
\begin{itemize}
\item Every pivotal rank-five fusion category admits a spherical structure.
\item If its universal grading is trivial, every pivotal structure is spherical.
\end{itemize}
The proof now fixes the pivotal structure when discussing its slope.
In the nontrivially graded case it invokes weak integrality and the
canonical spherical structure from pseudo-unitarity. The pointed example
is included to make the distinction explicit. Theorem~A is unchanged.

\subsection{The odd-multiplicity indicator rule: M16}

\textbf{Location: Section 17, mathematical content of the finite certificates.}
The original sentence asserted that occurrence with odd multiplicity in
$X_j\otimes X_j^*$ forces $\nu_2(X_i)=+1$. It omitted the hypothesis
$X_i\simeq X_i^*$.

There is a counterexample within the seventeen-ring list itself. In
$\operatorname{Rep}(C_7\rtimes C_3)$, let $g$ be a nontrivial linear
character and let $V$ be the degree-three representation induced from
the character orbit $\{1,2,4\}$ of $C_7$. Then
\[
 V\otimes V^*\cong\mathbf1\oplus g\oplus g^2\oplus V\oplus V^*.
\]
The differences of elements of $\{1,2,4\}$ are zero three times and each
nonzero residue modulo seven once, which verifies the decomposition.
Here $V\not\simeq V^*$ and $\nu_2(V)=0$, although its multiplicity is one.

The repaired rule assigns zero to non-self-dual simples. For self-dual
simples, odd multiplicity in the canonically symmetric object
$X_j\otimes X_j^*$ excludes indicator $-1$, and hence forces $+1$.
Otherwise both signs are retained. This is already exactly the logic in
\path{criteria/higher_traces.py}, function \texttt{indicator\_choices}:
the duality test precedes the odd-multiplicity test. The error was in the
written description; no witness, exclusion or count changes.

\subsection{The conductor hypothesis in the cubic depth proposition: M14}

\textbf{Location: Proposition 15.9, labelled
\texttt{prop:cubic-four-depth}.}
The proof uses Theorem 15.5, whose dyadic conclusions require conductor
$F\in\{7,9\}$. In particular, it uses the unramified inert cubic
completion at two. The proposition's statement did not repeat this
restriction. The revised statement explicitly assumes $F\in\{7,9\}$.
Every subsequent application already lies in these two conductors, so
the final cubic classification and the sharp bounds are unaffected.

\section{Proof dependencies and explanations made explicit}

\subsection{The dimensional grading input: M02, M18}

\textbf{Location: Proposition 3.2 and bibliography.}
The equality of the simple-dimension field and $\mathbb Q(D)$ now cites
Gannon--Schopieray, Proposition 1.8, and explains its use. With
$K_1=\mathbb Q(d_i:i)$ and $K_0=\mathbb Q(D)$, the fusion ring has a
faithful $\operatorname{Gal}(K_1/K_0)$-grading. Its group is a quotient
of the universal grading group, so trivial universal grading forces
$K_1=K_0$. No integrality of the $d_i$ is assumed. A primary preprint
link, \href{https://arxiv.org/abs/1912.12260}{arXiv:1912.12260}, was added
to the bibliography alongside the existing journal reference.

\subsection{The simultaneous choice of embedding: M04}

\textbf{Location: end of Section 3.}
Several descent arguments choose $E$ maximal among its conjugates while
using simultaneous minima of $1/E$, $1/D$ and $E/D$. The new paragraph
explains why these choices are compatible.

Transporting the category by a Galois automorphism changes its spherical
data and intrinsic dimension $E$, but preserves the integer fusion rules
and their distinguished positive Frobenius--Perron row. After transporting,
we choose the positive lifts again. Independently, the spectral-radius
bound gives $|\tau(d_i)|\le d_i$, hence $\tau(D)\le D$. Thus $1/D$
and $1/E$ are minimal at the chosen embedding, and the simultaneous-column
lemma, applied again to the transported category, makes $E/D$ minimal
there too. This supplies an implicit justification; it changes no
arithmetic range.

\subsection{A self-contained ordinary Fourier bound: M05, M06, M12, M13, M15}

\textbf{Locations: new Proposition 4.4, introduction, and the quadratic,
cubic and quartic applications.}
The new proposition proves the ordinary-center estimate directly from
Lemma 4.1. Suppose that the category has positive spherical dimensions,
commutative Grothendieck ring, and
$K=\mathbb Q(d_i:i)=\mathbb Q(D)$. For $d=[K:\mathbb Q]$ and a
Dirichlet character $\eta$ of conductor $c$, it gives
\[
 \frac1d\sum_\rho\left|\sum_{\sigma\in\operatorname{Gal}(K/\mathbb Q)}
 \frac{\eta(\sigma)}{\sigma(f_\rho)}\right|^2\le\frac1c,
 \qquad
 \left|\sum_\sigma\frac{\eta(\sigma)}{\sigma(D)}\right|
 \le c^{-1/2}.
\]
The proof uses the center of the category itself, normalized by $D$.
The $d$ vacuum-orbit labels have twist one, and their entries at the
unit-induction constituents are $1/\sigma(f_\rho)$. Applying the
unipotent projection bound to their normalized Fourier vector gives
the first inequality. The $d$ dimension-orbit terms have equal absolute
Fourier values, giving the second. The proof states why $f_\rho\in K$
and specifies the cyclotomic inflation of $\eta$.

The introduction now describes this common proof. The quadratic profile
reduction and quartic field list cite the internal proposition. The
repeated ordinary Fourier derivation in Proposition 15.3 has been
replaced by its direct application. Schopieray's attribution is retained.
The full text of the recent dimension-block preprint was not accessible
during this review; the new derivation makes the estimate used here
independent of that access limitation. Its underlying congruence and
modular-character inputs remain explicitly cited.

\subsection{The ramified version of the local-order argument: M09}

\textbf{Location: proof of Lemma 13.2.}
The proof now explicitly records the extension already used in the
quadratic profile argument. For a local fiber product
$O\times_{\mathbb F_\ell}B$ with a perfect associative pairing, the
first factor may be any discrete valuation ring finite over
$\mathbb Z_\ell$ with residue field $\mathbb F_\ell$, including a
ramified quadratic integer ring.

If $B$ were nonnormal, choose $y$ outside it with
$\mathfrak m_B y\subset\mathfrak m_B$. The classes $(\ell,0)$ and
$(0,\ell y)$ give two independent socle vectors modulo $\ell$,
contrary to perfect duality. Thus $B$ is a valuation ring. In the
codegree-five quadratic application, the trace dual for
$\Tr(xy/5)$ is the ramified prime ideal, so the gluing quotient really
is $\mathbb F_5$. This confirms and explains the existing argument;
it does not replace it by an assumption.

\subsection{The rotation underlying the local indicator congruence: M10}

\textbf{Location: proof of Lemma 13.4.}
The grouping step now names its operator. For $M=m\ell^n$, put
$R_n=E^{(M,1)}_{(i,i^*),V}$. Ng--Schauenburg,
Lemma 2.5 and Proposition 2.8(ii), identify the two indicators with
\[
 \Tr(R_n^{\ell^n})\quad\hbox{and}\quad
 \Tr(R_n^{\ell^{n-b}}).
\]
Their Lemma 2.7 gives $R_n^M=1$ because the insertion has twist one.
Their Proposition 5.7 puts all power traces in $\mathbb Q(\zeta_N)$.
The ensuing Galois-orbit multiplicity argument and divisibility
$\ell^{n-2b+1}\mid\ell^{t-b}$ with $t-b\ge n-2b+1$ are unchanged.
The exact source identities were checked in
\href{https://arxiv.org/abs/0806.2493}{arXiv:0806.2493}.

\subsection{Reference scope and small notation repairs: M07, M08, M11, M17}

In Corollary 10.1, the application of the near-group classifications now
explicitly takes place in the positive spherical conjugate already
provided by Theorem A. Thus its use of Meng--Yu does not require another
pivotality or positivity assertion. The full primary Meng--Yu manuscript
was checked, as were the twist-order statement in Schopieray--Wang--Yu,
Proposition 2.2(2), and the relevant unitary realizations.

In the nonreal-character proposition in Section 12.2, the display now
writes $f_\chi=h$ explicitly. In Proposition 14.2, the dimension-vector
display explicitly says that its sums range over the four nonunit
simples; this prevents inadvertently counting the unit twice in
$A=4+\sum a_i^2$. The Schopieray near-group journal reference was
completed with issue 4 and pages 1617--1641, as verified in the primary
bibliography. These are notation and reference corrections, not new
classification arguments.

\section{Program corrections and the new reproducible audit}

\subsection{Outward rounding of the generator bound: C01}

\textbf{File: \path{checks/small_simple.py}.}
The old expression \texttt{d0//SCALE+1} converted an approximate
embedding value into an asserted integer upper endpoint. The value was
computed from a root endpoint and rounded powers, with polynomial
coefficients of both signs. The expression did not itself include the
known approximation error.

The revised expression adds
\texttt{upper\_error = 15*sum(abs(v) for v in x)} before rounding
upwards. Each scaled power has error less than 15 by the existing
rational enclosures; the triangle inequality therefore encloses the
whole polynomial. This makes the subsequent finite range for
$N_{XX}^X$ rigorously outward. The three generator cases retain exactly
931, 1158 and 856 dimension candidates, respectively, and no final
survivor. No omitted actual candidate was found in the original run.

\subsection{Refusing to verify with assertions disabled: C02}

The top-level \path{verify.py} already rejected optimized Python, but
some standalone scripts could print a success message after
\texttt{python -O} had removed their assertions. An immediate
\texttt{if not \_\_debug\_\_} guard now rejects that invocation in
\path{final_census/verify.py} and these eleven structural programs:
\begin{quote}\small\raggedright
\path{biquad_sieve.py}, \path{cyclic_dimension.py},
\path{quadratic_profiles.py}, \path{quartic_five_twists.py},
\path{small_simple.py}, \path{verify15.py},
\path{verify_biquartic.py}, \path{verify_cyclic_quartic.py},
\path{verify_quadratic_120.py}, \path{verify_quartic_moments.py},
\path{verify_quintic_trace.py}.
\end{quote}
All thirteen guarded entry points, including the unchanged top-level
one, were tested with \texttt{-O} and rejected execution. The other seven
structural programs use explicit exceptions for their decisive checks.
No mathematical condition was added or removed by these guards.

\subsection{Stale documentation inside the induction routine: C03}

In \path{final_census/criteria/induction.py}, an undefined
\texttt{[DP]} attribution was replaced by the actual tests implemented.
The restriction-column docstring was corrected from the rank-four
$(1,v_1,v_2,v_3)$ to the rank-independent
$(1,v_1,\ldots,v_{n-1})$. Neither edit changes the algorithm.

\subsection{Restoring a missing audit program: C04}

A generator header referred to \path{verify_independent_census.py},
but the supplied archive did not contain it. The new
\path{final_census/verify_independent_census.py} now provides a
reproducible cross-check. It reconstructs the symmetry orbits and
associativity equations from multiplication tensors, checks the stored
and C++ equations symbolically, and verifies that every equation used
in linear completion is necessary. It replays the bounded census with
the nonlinear shortcuts bypassed and compares the resulting sets.

The slower replay shares the existing fraction-free linear-completion
routine and the basis parametrization. Its documentation states this;
it is not presented as a wholly independent implementation or a formal
verification. It also constructs the neutral $\mathrm{SU}(3)_4$ tensor
from its level-four Pieri graph and compares it with the stored model.
No witness records are needed for this optional audit.

\section{Verification results and mathematical checks}

Both complete runs reproduce the same finite census:
\begin{center}
\begin{tabular}{@{}rrrrr@{}}
\toprule
$\mu$ & Generated & Within both bounds & Excluded & Realized\\\midrule
1&16&16&6&10\\
2&37&30&27&3\\
3&82&46&44&2\\
4&134&22&20&2\\\midrule
Total&269&114&97&17\\\bottomrule
\end{tabular}
\end{center}
The 97 exclusions consist of 76 formal-codegree $d$-number failures,
two subcategory-divisibility failures, three central-dimension norm
failures, two dimension-field failures, ten induction/second-moment
failures and four higher-indicator certificates. These are exclusive
certificate categories, not assertions that each ring fails only one
criterion. All seventeen positive tensors match their models under
unit-preserving permutations.

The optional audit gives 220 self-dual rings, 47 rings with one
nontrivial dual pair and two rings with two such pairs. The respective
associativity equation counts are 21, 12 and 13, and the checked
linear-completion equation counts are 15, 11 and 12. The optimized and
slower censuses agree exactly. The independent Pieri construction
agrees with the $\operatorname{Ad}(\mathrm{SU}(3)_4)$ registry entry.

All eighteen structural checks passed. Their main outputs, useful for
matching the supplied transcript to the mathematical reductions, are:
\begin{center}\small
\begin{tabularx}{\textwidth}{@{}>{\raggedright\arraybackslash}p{.23\textwidth}>{\raggedright\arraybackslash}X@{}}
\toprule
Calculation & Reproduced result\\\midrule
Quintic descent & 34 ratios; 15,396 trace vectors; 39 signed passes;
four joint passes, all Galois conjugate.\\
Quadratic descent & Successive counts 331, 50, 13, 22, 1; the one final
tuple is excluded by the written complementary-unit argument.\\
Biquadratic descent & 2,097 fields, 52 after the field bounds;
30 ratios, all excluded by integral-moment Fourier tests.\\
Cyclic quartic descent & 105 fields, 221 field-polynomial profiles,
32 ratio profiles, ten norm candidates, no Fourier-compatible tuple.\\
Quadratic profiles & 1,299 rational and 614 mixed finite profiles;
162 rational and six mixed limiting profiles. The five remaining
dimension-vector cases fail.\\
Cubic dimension & Only $(r,t,F,A,sN,N)=(7,7,7,42,245,343)$ remains;
the polar conductor-seven profile has no candidate.\\
Quartic generators & For $(F,f)=(15,4),(15,5),(20,5)$:
931, 1,158, 856 dimensions; 20, 3, 23 norm candidates;
2, 0, 3 before Fourier; no final survivor.\\
Exchanged-pair moments & Ten root columns and six target tuples at
$f=20$; 25 columns and 180 target tuples at $f=5$; all excluded.\\\bottomrule
\end{tabularx}
\end{center}
The remaining checks reproduce the quartic field table, all ten moment
bounds, the generator thresholds 125, 155 and 153, the codegree-sixteen
inequality, and the biquadratic and cyclic dimension exclusions. The
full unabridged output is in \path{verification.txt}.

Several potentially serious points were checked without requiring a
correction: the commutativity argument for the signed order; the
unconditional rank-three descent input; normalization by $2E$ in the
sphericalization; the use of a unitary modular representation without
assuming a unitary category; both Hecke estimates; the signed unit
argument; the local socle obstruction; the Galois-orbit indicator
congruence; the Morita construction; the ADE generator reduction; and
the quintic trace lattice. The finite dimension test was also checked:
its Casimir matrix is real symmetric, and nonnegative coefficients of
the shifted real-rooted characteristic polynomial are equivalent to
the required upper bound on its largest eigenvalue.

The unitary endpoint construction was checked directly against Izumi,
Lemma 13.1 and Example 13.10. The relevant symmetry is an actual
involution, fixes the unit intertwiner, and has trace zero on the
six-dimensional noninvertible multiplicity space. It therefore gives
the stated self-dual equivariantization, multiplicities three and
$D=48+24\sqrt3$. The categorical realization and the exact tensor
comparison serve separate purposes; matching arrays alone is not a
proof of unitarity.

\section{Complete package changes and reproduction}\label{sec:package}

Beyond M01--M18 and C01--C04, the delivery changes are as follows:
\begin{enumerate}
\item The current manuscript PDF was rebuilt from the revised source.
The title, author fields, seventeen-model list, principal conclusions,
and sharp constants were preserved. The extra proposition is numbered
4.4; later sections retain their local theorem numbering.
\item This new letter is supplied as \path{Revision_Letter.tex} and
\path{Revision_Letter.pdf}. The old review source and PDF remain unchanged.
\item The root README now identifies the current letter, distinguishes
the previous report, explains the corrected assertion-mode behavior,
and documents the optional census audit. The census README describes
that audit and its shared implementation components.
\item \path{verification.txt} was refreshed with the complete successful
run of the revised programs. \path{independent_verification.txt} records
the additional audit results and the assertion-mode rejection checks.
The census receipt is consistent with the new run; its numerical
content is unchanged.
\item \path{revision.patch} records every edit to an existing text file,
excluding the regenerated checksum manifest. New files are individually
identified above. The manifest \path{SHA256SUMS.txt} was regenerated for
the final distribution. Temporary build files, compiled helpers and
Python cache files are excluded.
\end{enumerate}

No stored witness tensor, model tensor, census-generator equation,
mathematical rejection criterion or dependency requirement was changed.
The correction to the endpoint enclosure is the only change to a
numerical search bound; its observed results are unchanged.

From the extracted source directory, run:
\begin{verbatim}
python3 verify.py
python3 final_census/verify_independent_census.py
\end{verbatim}
The second command is optional and performs the additional census/model
audit. Requirements remain Python with SymPy and NumPy, and a C++17
compiler with OpenMP. The supplied programs use exact arithmetic for
acceptance and exclusion. Normal verification requires no network access
or numerical feasibility solver. The LaTeX sources compile with a
standard installation of pdfLaTeX and the packages in their preambles.

For exact identification, the SHA-256 of the incoming ZIP is
\begin{quote}\small\raggedright\path{c89ab1ffb252cacb2d135dbdcde5bb1d61803bacdb9eccdd21f810ad625ab63e}.
\end{quote}
The SHA-256 of its original \path{rank5.tex} is
\begin{quote}\small\raggedright\path{59acf340df9863608c6e0e359dc375d055b5132bc102dd9974f9a20d093426d2}.
\end{quote}
These identify the baseline against which every change in this letter
and the patch is measured.

\end{document}
